% =====================================================================
% CAPÍTULO 2: METODOLOGÍA DE INVESTIGACIÓN
% =====================================================================
\chapter{Metodología de Investigación}
\label{ch:metodologia}

\section{Enfoque de Investigación}

La presente investigación adopta un enfoque \textbf{mixto} (cualitativo y
cuantitativo) con alcance \textbf{descriptivo, explicativo y propositivo}. Se
eligió este enfoque porque el problema del \ac{DIC} combina dimensiones humanas
(percepciones, comunicación, conflictos) con dimensiones medibles (tiempos de
respuesta, tasas de respuesta, frecuencias de quejas, disponibilidad horaria).

\begin{table}[H]
\centering
\caption{Componentes del enfoque mixto}
\label{tab:enfoque-mixto}
\begin{tabularx}{\textwidth}{@{}l X X@{}}
\toprule
\rowcolor{azulClaro}
\textbf{Dimensión} & \textbf{Enfoque cualitativo} & \textbf{Enfoque cuantitativo} \\
\midrule
Propósito & Comprender el problema y sus causas & Medir magnitud, frecuencia e impacto \\
Técnicas & Entrevistas, observación, análisis documental & Encuestas, registro de tiempos, analítica \\
Instrumentos & Guion de entrevista, ficha de observación & Cuestionario estructurado, \ac{ETL} de logs \\
Muestra & 8--12 informantes clave & Muestreo probabilístico por estrato \\
Análisis & Codificación temática & Estadística descriptiva e inferencial \\
Resultado & Marco interpretativo del problema & Indicadores y modelos predictivos \\
\bottomrule
\end{tabularx}
\end{table}

\section{Tipo y Diseño de Investigación}

\subsection{Tipo}
\begin{itemize}[leftmargin=1.4cm, itemsep=3pt]
  \item \textbf{Investigación aplicada:} busca resolver un problema concreto del
        \ac{DIC} mediante una solución tecnológica documentada.
  \item \textbf{Estudio de caso:} el \ac{DIC} se analiza como unidad integral,
        considerando sus procesos, actores y datos.
  \item \textbf{Investigación--acción:} se iteran ciclos de diagnóstico, diseño,
        validación y refinamiento con participación de los actores.
\end{itemize}

\subsection{Diseño}
Se emplea un diseño \textbf{transversal} para el diagnóstico (una medición
completa del estado actual) y \textbf{longitudinal} para la validación (mediciones
del antes y después de la implementación de los módulos).

\begin{figure}[H]
\centering
\begin{tikzpicture}[node distance=0.55cm and 1.1cm, every node/.style={font=\scriptsize}]

\node[terminal] (inicio) {Inicio};
\node[flujo, below=of inicio] (f1) {Diagnóstico del\\estado actual};
\node[flujo, below=of f1] (f2) {Recolección cualitativa\\y cuantitativa};
\node[flujo, below=of f2] (f3) {Análisis y\\codificación};
\node[decision, below=of f3, yshift=-0.35cm] (d1) {¿Problema\\validado?};
\node[flujo, left=of d1, xshift=-0.9cm] (f4) {Ampliar\\muestra};
\node[flujo, below=of d1, yshift=-0.5cm] (f5) {Diseño de\\la solución};
\node[flujo, below=of f5] (f6) {Prototipo y\\validación};
\node[decision, below=of f6, yshift=-0.35cm] (d2) {¿Hipótesis\\aceptada?};
\node[flujo, right=of d2, xshift=0.9cm] (f7) {Refinar\\iterativamente};
\node[terminal, below=of d2, yshift=-0.5cm] (fin) {Documentación\\final};

\draw[flecha] (inicio) -- (f1);
\draw[flecha] (f1) -- (f2);
\draw[flecha] (f2) -- (f3);
\draw[flecha] (f3) -- (d1);
\draw[flecha] (d1) -- node[above, font=\tiny]{No} (f4);
\draw[flecha] (f4) |- (f2);
\draw[flecha] (d1) -- node[right, font=\tiny]{Sí} (f5);
\draw[flecha] (f5) -- (f6);
\draw[flecha] (f6) -- (d2);
\draw[flecha] (d2) -- node[above, font=\tiny]{No} (f7);
\draw[flecha] (f7) |- (f5);
\draw[flecha] (d2) -- node[right, font=\tiny]{Sí} (fin);

\end{tikzpicture}
\caption{Ciclo metodológico de investigación--acción aplicado al \ac{DIC}}
\label{fig:ciclo-metodologico}
\end{figure}

\section{Fases de la Metodología}

La metodología se estructura en ocho fases, cada una con entradas, actividades y
entregables explícitos.

\begin{table}[H]
\centering
\caption{Fases de la metodología de investigación}
\label{tab:fases}
\small
\begin{tabularx}{\textwidth}{@{}c l X l@{}}
\toprule
\rowcolor{azulClaro}
\textbf{Fase} & \textbf{Nombre} & \textbf{Actividades principales} & \textbf{Duración} \\
\midrule
F1 & Definición del problema &
    Identificación de síntomas, delimitación del alcance, formulación de preguntas
    e hipótesis. & 2 sem. \\
F2 & Revisión de literatura &
    Búsqueda de antecedentes en gestión académica, analítica educativa, sistemas
    de quejas y evaluación docente. & 3 sem. \\
F3 & Diseño metodológico &
    Selección de técnicas, definición de población y muestra, construcción de
    instrumentos. & 2 sem. \\
F4 & Recolección de datos &
    Entrevistas, encuestas, observación de procesos, extracción de datos
    históricos del \ac{DIC}. & 4 sem. \\
F5 & Análisis de datos &
    Transcripción y codificación, estadística descriptiva, análisis de
    correlaciones, preparación del dataset. & 3 sem. \\
F6 & Diseño de la solución &
    Arquitectura, modelado \ac{UML}, especificación de requerimientos y del
    algoritmo de soluciones rápidas. & 4 sem. \\
F7 & Validación &
    Validación de instrumentos, validación del modelo con expertos, prueba piloto
    del prototipo. & 3 sem. \\
F8 & Documentación y difusión &
    Redacción del informe, tableros de resultados, socialización con la comunidad
    académica. & 3 sem. \\
\bottomrule
\end{tabularx}
\end{table}

\section{Población y Muestra}

\subsection{Población}
La población objeto de estudio está conformada por los integrantes del \ac{DIC}:

\begin{itemize}[leftmargin=1.4cm, itemsep=3pt]
  \item \textbf{Estudiantes:} matriculados en materias del departamento en el
        periodo académico vigente.
  \item \textbf{Docentes:} adscritos al departamento.
  \item \textbf{Personal de secretaría:} responsable de la atención de
        solicitudes, quejas y sugerencias.
  \item \textbf{Directivos:} dirección y coordinación académica del departamento.
\end{itemize}

\subsection{Tamaño de muestra}
Para la encuesta a estudiantes se utiliza muestreo probabilístico estratificado
por semestre, con un nivel de confianza del $95\%$ y un margen de error del $5\%$:

\begin{equation}
n = \frac{N \cdot Z^{2} \cdot p \cdot q}{e^{2}(N-1) + Z^{2} \cdot p \cdot q}
\label{eq:muestra}
\end{equation}

donde $N$ es el tamaño de la población, $Z = 1.96$, $p = q = 0.5$ y
$e = 0.05$. Para el resto de actores se emplea \textbf{muestreo no probabilístico
por conveniencia}, dado el tamaño reducido de la población (docentes, secretaría
y directivos).

\begin{table}[H]
\centering
\caption{Distribución de la muestra por actor}
\label{tab:muestra}
\begin{tabularx}{\textwidth}{@{}l c c X@{}}
\toprule
\rowcolor{verdeClaro}
\textbf{Actor} & \textbf{Población} & \textbf{Muestra} & \textbf{Técnica principal} \\
\midrule
Estudiantes & $N$ (variable) & $n$ calculado & Encuesta estructurada \\
Docentes & 20--40 & 12--20 & Entrevista semiestructurada \\
Secretaría & 2--4 & 2--4 & Entrevista + observación de procesos \\
Directivos & 2--3 & 2--3 & Entrevista en profundidad \\
Datos históricos & Registros del \ac{DIC} & Censo & Extracción y \ac{ETL} \\
\bottomrule
\end{tabularx}
\end{table}

\section{Técnicas e Instrumentos de Recolección}

\begin{table}[H]
\centering
\caption{Técnicas e instrumentos de recolección de datos}
\label{tab:instrumentos}
\small
\begin{tabularx}{\textwidth}{@{}p{3.4cm} p{4.9cm} X@{}}
\toprule
\rowcolor{naranjaClaro}
\textbf{Técnica} & \textbf{Instrumento} & \textbf{Propósito} \\
\midrule
Encuesta & Cuestionario de 25 ítems (escala Likert 1--5) &
    Medir percepción de tiempos, satisfacción con comunicación y utilidad percibida \\
Entrevista semiestructurada & Guion de 15 preguntas &
    Comprender causas del retraso, sobrecarga operativa y expectativas \\
Observación no participante & Ficha de registro de procesos &
    Cronometrar el flujo real de una solicitud de cupo y de una queja \\
Análisis documental & Ficha de revisión normativa &
    Identificar reglas de negocio y criterios de asignación vigentes \\
Cuestionario de horarios & Formulario de disponibilidad y preferencia &
    Capturar datos para el modelo probabilístico de horarios \\
Extracción de datos & Consultas sobre registros históricos &
    Construir el dataset para analítica de datos \\
Grupo focal & Guion de discusión (6--8 estudiantes) &
    Validar hallazgos y priorizar funcionalidades \\
\bottomrule
\end{tabularx}
\end{table}

\subsection{Escala de medición}
Se emplea una escala Likert de cinco niveles para las variables perceptuales:

\begin{table}[H]
\centering
\caption{Escala Likert utilizada en los cuestionarios}
\label{tab:likert}
\begin{tabular}{@{}c l@{}}
\toprule
\rowcolor{moradoClaro}
\textbf{Valor} & \textbf{Significado} \\
\midrule
5 & Totalmente de acuerdo / Muy satisfecho \\
4 & De acuerdo / Satisfecho \\
3 & Neutral / Indiferente \\
2 & En desacuerdo / Insatisfecho \\
1 & Totalmente en desacuerdo / Muy insatisfecho \\
\bottomrule
\end{tabular}
\end{table}

\subsection{Operacionalización de variables}
\begin{table}[H]
\centering
\caption{Operacionalización de variables}
\label{tab:variables}
\small
\begin{tabularx}{\textwidth}{@{}l l l X@{}}
\toprule
\rowcolor{azulClaro}
\textbf{Variable} & \textbf{Tipo} & \textbf{Indicador} & \textbf{Instrumento} \\
\midrule
Tiempo de respuesta & Cuantitativa continua & Horas promedio & Observación, logs \\
Tasa de respuesta docente & Cuantitativa discreta & Porcentaje & Registro del sistema \\
Satisfacción comunicacional & Cualitativa ordinal & Índice 1--5 & Encuesta \\
Frecuencia de quejas & Cuantitativa discreta & N.º por periodo & Registro del sistema \\
Preferencia de horario & Cualitativa nominal & Probabilidad $P(h_i)$ & Cuestionario de horarios \\
Desempeño docente & Cuantitativa continua & Puntaje 0--5 & Evaluación docente \\
Efectividad secretaría & Cuantitativa continua & $\%$ casos resueltos & Tablero verificación \\
\bottomrule
\end{tabularx}
\end{table}

\section{Metodología CRISP-DM para la Analítica de Datos}

Para el componente de analítica de datos y el algoritmo de soluciones rápidas se
adopta el estándar \ac{CRISP-DM}, que organiza el trabajo de datos en seis fases
iterativas.

\begin{figure}[H]
\centering
\begin{tikzpicture}[
  fase/.style={
    draw=azulOscuro, fill=azulClaro, thick, rounded corners=6pt,
    text=azulOscuro, align=center, font=\small\bfseries,
    minimum width=3.2cm, minimum height=1.1cm, inner sep=4pt
  },
  node distance=0.7cm and 0.9cm
]

\node[fase] (f1) {1. Comprensión\\del negocio};
\node[fase, right=of f1] (f2) {2. Comprensión\\de los datos};
\node[fase, right=of f2] (f3) {3. Preparación\\de los datos};
\node[fase, below=of f3] (f4) {4. Modelado};
\node[fase, left=of f4] (f5) {5. Evaluación};
\node[fase, left=of f5] (f6) {6. Despliegue};

\draw[flecha] (f1) -- (f2);
\draw[flecha] (f2) -- (f3);
\draw[flecha] (f3) -- (f4);
\draw[flecha] (f4) -- (f5);
\draw[flecha] (f5) -- (f6);
\draw[flecha] (f6) -- (f1);
\draw[-{Stealth[length=2mm,width=1.4mm]}, thick, dashed, color=moradoUC, bend right=35]
      (f4) to node[below, font=\tiny]{iterar} (f3);
\draw[-{Stealth[length=2mm,width=1.4mm]}, thick, dashed, color=moradoUC, bend left=35]
      (f5) to node[above, font=\tiny]{refinar} (f2);

\end{tikzpicture}
\caption{Ciclo \ac{CRISP-DM} adaptado a la analítica del \ac{SGDIC}}
\label{fig:crispdm}
\end{figure}

\begin{table}[H]
\centering
\caption{Aplicación de \ac{CRISP-DM} al \ac{SGDIC}}
\label{tab:crispdm}
\small
\begin{tabularx}{\textwidth}{@{}c l X@{}}
\toprule
\rowcolor{verdeClaro}
\textbf{Fase} & \textbf{Nombre} & \textbf{Aplicación concreta en el \ac{DIC}} \\
\midrule
1 & Comprensión del negocio &
    Objetivos: reducir tiempos de respuesta, asegurar evaluación docente
    oportuna, priorizar quejas. Definición de \ac{KPI}. \\
2 & Comprensión de los datos &
    Fuentes: registros de cupos, mensajes, faltas, evaluaciones, quejas,
    cuestionarios de horarios. \\
3 & Preparación de los datos &
    Limpieza, normalización de fechas y categorías, anonimización de \ac{PII},
    construcción de variables derivadas. \\
4 & Modelado &
    Modelos de clasificación de quejas, predicción de demanda de cupos,
    probabilidad de preferencia de horario, detección de anomalías. \\
5 & Evaluación &
    Validación cruzada, matrices de confusión, comparación contra línea base,
    revisión con expertos del departamento. \\
6 & Despliegue &
    Integración de las recomendaciones en el tablero de verificación y en el
    motor de soluciones rápidas. \\
\bottomrule
\end{tabularx}
\end{table}

\section{Validez y Confiabilidad}

\subsection{Validez de contenido}
Los instrumentos se someten a \textbf{juicio de expertos} (tres docentes con
experiencia en metodología de investigación y dos directivos del \ac{DIC}), con
una rúbrica de evaluación de claridad, pertinencia y suficiencia.

\subsection{Validez de constructo}
Se aplica \textbf{análisis factorial exploratorio} para verificar que los ítems
del cuestionario agrupan las dimensiones previstas: comunicación, oportunidad,
trazabilidad y satisfacción.

\subsection{Confiabilidad}
La consistencia interna se mide con el coeficiente \textbf{alfa de Cronbach}:

\begin{equation}
\alpha = \frac{k}{k-1}\left(1 - \frac{\sum_{i=1}^{k} \sigma^{2}_{i}}{\sigma^{2}_{T}}\right)
\label{eq:cronbach}
\end{equation}

donde $k$ es el número de ítems, $\sigma^{2}_{i}$ la varianza de cada ítem y
$\sigma^{2}_{T}$ la varianza total. Se considera aceptable un valor
$\alpha \geq 0.70$.

\subsection{Validez de la analítica}
Para los modelos de datos se emplea \textbf{validación cruzada de $k$ particiones}
($k=10$) y se reportan exactitud, precisión, exhaustividad y $F_1$:

\begin{equation}
F_1 = 2 \cdot \frac{\text{Precisión} \cdot \text{Exhaustividad}}
{\text{Precisión} + \text{Exhaustividad}}
\label{eq:f1}
\end{equation}

\section{Consideraciones Éticas}

\begin{cajaNota}[Principios éticos del estudio]
\begin{itemize}[leftmargin=1.2cm, itemsep=3pt]
  \item \textbf{Consentimiento informado:} todos los participantes conocen el
        propósito del estudio y autorizan su participación.
  \item \textbf{Confidencialidad:} las respuestas de estudiantes sobre docentes
        se anonimizan; ningún docente conoce la identidad de quien lo califica.
  \item \textbf{No represalia:} se garantiza que reportar una queja o una
        evaluación negativa no genere consecuencias académicas adversas.
  \item \textbf{Minimización de datos:} solo se recolectan los datos necesarios
        para cada objetivo declarado.
  \item \textbf{Transparencia algorítmica:} los criterios del algoritmo de
        priorización son públicos y auditables por la comunidad académica.
  \item \textbf{Revisión institucional:} el protocolo se somete al comité de
        ética o instancia equivalente de la facultad.
\end{itemize}
\end{cajaNota}

\section{Análisis de Datos Previsto}

\begin{table}[H]
\centering
\caption{Técnicas de análisis por tipo de dato}
\label{tab:analisis}
\small
\begin{tabularx}{\textwidth}{@{}p{3.4cm} p{5.4cm} X@{}}
\toprule
\rowcolor{azulClaro}
\textbf{Tipo de dato} & \textbf{Técnica} & \textbf{Resultado esperado} \\
\midrule
Respuestas Likert & Estadística descriptiva (media, mediana, desviación) &
    Índice de satisfacción por dimensión \\
Comparación antes/después & Prueba $t$ para muestras relacionadas &
    Significancia del cambio en tiempos de respuesta \\
Datos categóricos de quejas & Tablas de contingencia, $\chi^{2}$ &
    Asociación entre tipo de queja y área responsable \\
Series temporales de cupos & Suavizado exponencial, ARIMA &
    Pronóstico de demanda por materia \\
Preferencias de horario & Frecuencias y probabilidad condicional &
    $P(\text{horario} \mid \text{perfil del estudiante})$ \\
Texto libre de quejas & PLN (bolsa de palabras) &
    Categorización automática y temas recurrentes \\
Registros de evaluación & Tendencia y control estadístico &
    Detección de docentes con desempeño atípico \\
\bottomrule
\end{tabularx}
\end{table}

\section{Limitaciones Metodológicas}

\begin{itemize}[leftmargin=1.4cm, itemsep=3pt]
  \item \textbf{Sesgo de autoselección:} quienes responden encuestas pueden tener
        opiniones más extremas que quienes no participan.
  \item \textbf{Deseabilidad social:} los estudiantes podrían suavizar sus
        evaluaciones por temor a represalias, mitigado con anonimato estricto.
  \item \textbf{Datos históricos incompletos:} los registros previos pueden ser
        fragmentarios, lo que obliga a imputar o descartar observaciones.
  \item \textbf{Estacionalidad académica:} los patrones de quejas y demanda de
        cupos varían según el calendario académico, lo que exige comparar
        periodos equivalentes.
  \item \textbf{Generalización:} los hallazgos describen el \ac{DIC} y no se
        extrapolan automáticamente a otras unidades académicas.
\end{itemize}
